\documentclass[11pt]{article}
\usepackage[margin=1in]{geometry}
\usepackage{amsmath, amssymb}
\usepackage{booktabs}
\usepackage[hidelinks]{hyperref}

\newcommand{\one}{\mathbf{1}}
\newcommand{\bv}{\mathbf{v}}
\newcommand{\br}{\mathbf{r}}
\newcommand{\bw}{\mathbf{w}}
\newcommand{\E}{\mathbb{E}}
\newcommand{\Cov}{\mathrm{Cov}}
\newcommand{\Var}{\mathrm{Var}}

\title{The multi-epoch likelihood in \texttt{multiepoch\_nll\_per\_star}}
\author{memo for \texttt{single\_component\_comparisons\_mockobs.py}}
\date{15 November 2026}

\begin{document}
\maketitle

\section*{Summary}
Each star's epochs are modeled as one draw from a multivariate gaussian.
The shared-stream model and the independent-epoch model have the \emph{same diagonal} covariance entries.
They differ only in the off-diagonal entries: $\sigma_s^2$ in the shared model, $0$ in the independent one.
The code doesn't build or invert a covariance matrix per star.
Because the matrix is ``diagonal plus a constant'', its inverse and determinant have closed forms
(Sherman--Morrison and the matrix determinant lemma).
The log likelihood then reduces to a few sums over epochs, which vectorize over all stars at once.
This memo derives those closed forms.

\section{The model}
Take one star, observed at epochs $k = 1, \dots, n$ (the code uses $n=2$).
The measured line-of-sight velocity is
\begin{equation}
  v_k = \mu + s + b_k + e_k ,
\end{equation}
where
\begin{itemize}
  \item $\mu$ is the stream track (a parameter);
  \item $s \sim \mathcal{N}(0, \sigma_s^2)$ is the star's offset from the track, drawn \textbf{once} per star (no $k$ subscript);
  \item $b_k \sim \mathcal{N}(0, \sigma_b^2)$ is the binary orbital motion, drawn fresh at each epoch;
  \item $e_k \sim \mathcal{N}(0, \epsilon_k^2)$ is the measurement noise, drawn fresh at each epoch ($\epsilon_k$ = \texttt{v\_err}).
\end{itemize}
All of $s, b_1, \dots, b_n, e_1, \dots, e_n$ are mutually independent with mean zero.
In the independent-epoch model, $s$ is replaced by $s_k$, a fresh draw at each epoch.

Throughout, define
\begin{equation}
  d_k \equiv \epsilon_k^2 + \sigma_b^2, \qquad r_k \equiv v_k - \mu .
\end{equation}
Bold symbols are length-$n$ column vectors over epochs, and $\one = (1, \dots, 1)^T$.

\section{The covariance matrix}
$v_k$ is a sum of gaussians, so the vector $\bv$ is multivariate gaussian with mean $\mu\one$.
Its covariance matrix has entries $C_{kl} = \Cov(v_k, v_l) = \E[(v_k - \mu)(v_l - \mu)]$.
Write $v_k - \mu = s + b_k + e_k$ and multiply out:
\begin{equation}
  C_{kl} = \E\big[(s + b_k + e_k)(s + b_l + e_l)\big].
\end{equation}
For independent, mean-zero $X$ and $Y$, $\E[XY] = \E[X]\,\E[Y] = 0$.
So every cross term vanishes, and only terms that pair a variable with \emph{itself} survive:
\begin{itemize}
  \item $k = l$: $s\cdot s$, $b_k\cdot b_k$ and $e_k \cdot e_k$ survive, giving $C_{kk} = \sigma_s^2 + \sigma_b^2 + \epsilon_k^2 = \sigma_s^2 + d_k$.
  \item $k \neq l$: $b_k$ and $b_l$ are different draws, as are $e_k$ and $e_l$. Only $s\cdot s$ survives, giving $C_{kl} = \sigma_s^2$.
\end{itemize}
So
\begin{equation}
  C = D + \sigma_s^2\, \one\one^T, \qquad D = \mathrm{diag}(d_1, \dots, d_n).
  \label{eq:C}
\end{equation}
$\one\one^T$ is the $n\times n$ matrix of all ones.
For $n=2$:
\begin{equation}
  C_{\rm shared} = \begin{pmatrix} \sigma_s^2 + d_1 & \sigma_s^2 \\ \sigma_s^2 & \sigma_s^2 + d_2 \end{pmatrix},
  \qquad
  C_{\rm indep} = \begin{pmatrix} \sigma_s^2 + d_1 & 0 \\ 0 & \sigma_s^2 + d_2 \end{pmatrix}.
\end{equation}
In the independent model $\E[s_k s_l] = 0$ for $k \neq l$, so the off-diagonal terms are zero.

\section{The log likelihood}
The multivariate gaussian density gives
\begin{equation}
  -\ln p(\bv) = \frac{1}{2}\left[\, \br^T C^{-1} \br + \ln\det C + n \ln 2\pi \,\right],
  \qquad \br = \bv - \mu\one .
  \label{eq:mvn}
\end{equation}
Stars are independent, so the total is $-\ln\mathcal{L} = \sum_i -\ln p(\bv_i)$, with a different $C_i$ for each star
(the errors $\epsilon_{ik}$ differ).
Up to here this is just the textbook formula.
The rest of the memo evaluates $\br^T C^{-1}\br$ and $\ln\det C$ for the matrix in Eq.~\eqref{eq:C}
without numerically inverting anything.

\section{Evaluating it fast}
Define the vector $\bw = D^{-1}\one$, i.e.\ $w_k = 1/d_k$, and the two scalars
\begin{equation}
  A \equiv \one^T D^{-1} \one = \sum_k \frac{1}{d_k},
  \qquad
  B \equiv \one^T D^{-1} \br = \sum_k \frac{r_k}{d_k}.
\end{equation}
$A$ is the star's total precision, and $B$ is its precision-weighted summed residual.

\subsection{The inverse (Sherman--Morrison)}
For an invertible $D$ and vectors $\mathbf{u}$, $\mathbf{x}$,
$(D + \mathbf{u}\mathbf{x}^T)^{-1} = D^{-1} - \dfrac{D^{-1}\mathbf{u}\,\mathbf{x}^T D^{-1}}{1 + \mathbf{x}^T D^{-1}\mathbf{u}}$.
With $\mathbf{u} = \sigma_s^2\one$ and $\mathbf{x} = \one$ this gives
\begin{equation}
  C^{-1} = D^{-1} - \frac{\sigma_s^2}{1 + \sigma_s^2 A}\, \bw\bw^T .
  \label{eq:Cinv}
\end{equation}
\textit{Check}, by multiplying out and using $D\bw = \one$ and $\one^T\bw = A$:
\begin{align*}
  C\,C^{-1}
  &= \big(D + \sigma_s^2\one\one^T\big)\left(D^{-1} - \frac{\sigma_s^2}{1+\sigma_s^2 A}\bw\bw^T\right) \\
  &= I + \sigma_s^2\one\bw^T - \frac{\sigma_s^2}{1+\sigma_s^2 A}\one\bw^T - \frac{\sigma_s^4 A}{1+\sigma_s^2 A}\one\bw^T \\
  &= I + \sigma_s^2\one\bw^T\left[1 - \frac{1 + \sigma_s^2 A}{1+\sigma_s^2 A}\right] = I .
\end{align*}
Sandwiching Eq.~\eqref{eq:Cinv} between $\br^T$ and $\br$, and using $\bw^T\br = B$:
\begin{equation}
  \br^T C^{-1}\br = \sum_k \frac{r_k^2}{d_k} - \frac{\sigma_s^2 B^2}{1 + \sigma_s^2 A}.
  \label{eq:chi2}
\end{equation}

\subsection{The determinant (matrix determinant lemma)}
Factor out $D$:
\begin{equation*}
  \det C = \det\!\big(D(I + \sigma_s^2\bw\one^T)\big) = \det D \,\cdot\, \det(I + \sigma_s^2\bw\one^T).
\end{equation*}
For vectors $\mathbf{x}$, $\mathbf{y}$, the matrix $I + \mathbf{x}\mathbf{y}^T$ has eigenvalue $1 + \mathbf{y}^T\mathbf{x}$
(eigenvector $\mathbf{x}$).
Every vector orthogonal to $\mathbf{y}$ is also an eigenvector, with eigenvalue $1$.
So $\det(I + \mathbf{x}\mathbf{y}^T) = 1 + \mathbf{y}^T\mathbf{x}$.
Here $\mathbf{y}^T\mathbf{x} = \sigma_s^2\one^T\bw = \sigma_s^2 A$, so
\begin{equation}
  \ln\det C = \sum_k \ln d_k + \ln(1 + \sigma_s^2 A).
  \label{eq:logdet}
\end{equation}

\subsection{Result}
Putting Eqs.~\eqref{eq:chi2} and \eqref{eq:logdet} into Eq.~\eqref{eq:mvn}:
\begin{equation}
  \boxed{\;
  -\ln p(\bv) = \frac{1}{2}\left[
    \sum_k \frac{r_k^2}{d_k} - \frac{\sigma_s^2 B^2}{1 + \sigma_s^2 A}
    + \sum_k \ln d_k + \ln(1 + \sigma_s^2 A)
    + n\ln 2\pi \right]\;}
  \label{eq:result}
\end{equation}
with $r_k = v_k - \mu$, $d_k = \epsilon_k^2 + \sigma_b^2$, $A = \sum_k 1/d_k$ and $B = \sum_k r_k/d_k$.
This is exact, not an approximation, and holds for any $n$.
Checked numerically against \texttt{scipy.stats.multivariate\_normal} for $n = 1, 2, 3, 5$
(random $\mu, \sigma_s, \sigma_b, \epsilon_k$): the largest difference was $3\times10^{-12}$.

\paragraph{Why bother?}
For $n=2$ you could invert the $2\times 2$ matrix by hand.
The point is that Eq.~\eqref{eq:result} is only sums along the epoch axis.
So it vectorizes over all $N$ stars at once (arrays of shape $(N, n)$), for any $n$,
with no per-star matrix construction and no \texttt{np.linalg} calls.

\subsection{Mapping to the code}
\begin{center}
\begin{tabular}{ll}
\toprule
code & math \\
\midrule
\texttt{s2} & $\sigma_s^2$ \\
\texttt{d} & $d_k = \epsilon_k^2 + \sigma_b^2$ \\
\texttt{r} & $r_k = v_k - \mu$ \\
\texttt{A}, \texttt{B} & $A$, $B$ \\
\texttt{denom} & $1 + \sigma_s^2 A$ \\
\texttt{chi2} & $\br^T C^{-1}\br$, Eq.~\eqref{eq:chi2} \\
\texttt{logdet} & $\ln\det C$, Eq.~\eqref{eq:logdet} \\
return value & Eq.~\eqref{eq:result} \\
\bottomrule
\end{tabular}
\end{center}

\section{The two-epoch case written out}
For $n=2$ the inverse can be written directly, which shows what the off-diagonal term does.
With $C = \begin{pmatrix} a & c \\ c & b \end{pmatrix}$, $a = \sigma_s^2 + d_1$, $b = \sigma_s^2 + d_2$, $c = \sigma_s^2$:
\begin{equation*}
  \det C = ab - c^2 = d_1 d_2 + \sigma_s^2(d_1 + d_2) = d_1 d_2\,(1 + \sigma_s^2 A),
\end{equation*}
which agrees with Eq.~\eqref{eq:logdet}. Then
\begin{equation}
  \br^T C^{-1}\br = \frac{b r_1^2 - 2c\,r_1 r_2 + a r_2^2}{\det C}
  = \frac{1}{1 + \sigma_s^2 A}\left[\frac{r_1^2}{d_1} + \frac{r_2^2}{d_2} + \sigma_s^2\frac{(r_1 - r_2)^2}{d_1 d_2}\right],
  \label{eq:twoepoch}
\end{equation}
which agrees with Eq.~\eqref{eq:chi2}.
To check this, multiply Eq.~\eqref{eq:chi2} by $1 + \sigma_s^2A$ and expand $B^2$.

Eq.~\eqref{eq:twoepoch} shows what $\sigma_s$ is constrained by:
\begin{itemize}
  \item \textbf{A star off the track by the same amount at both epochs} ($r_1 = r_2 = r$, equal $d$):
        $\br^TC^{-1}\br = r^2/(\sigma_s^2 + d/2)$.
        This is the epoch mean judged against its own variance, $\sigma_s^2 + d/2$.
        A large $\sigma_s$ makes a consistent offset cheap.
  \item \textbf{Epochs that disagree} ($r_1 \neq r_2$): the $(r_1 - r_2)^2$ term is penalized with $d_k$ only.
        As $\sigma_s\to\infty$, $\br^TC^{-1}\br \to (r_1 - r_2)^2/(d_1 + d_2)$.
        So a disagreement between epochs can only be absorbed by $d_k$, i.e.\ by $\sigma_b$ or the noise, never by $\sigma_s$.
\end{itemize}
This is why the shared model can separate $\sigma_s$ from $\sigma_b$:
offsets that are consistent between epochs inform $\sigma_s$, and differences between epochs inform $\sigma_b$.

\section{Special cases}
\paragraph{Independent epochs (\texttt{shared\_stream=False}).}
$C = \mathrm{diag}(d_k + \sigma_s^2)$.
This is Eq.~\eqref{eq:C} with $d_k \to d_k + \sigma_s^2$ and the rank-one term set to zero.
The code does exactly that (\texttt{d, s2 = d + s2, 0.0}).
Eq.~\eqref{eq:result} then reduces to a product of 1D gaussians,
\begin{equation*}
  -\ln p = \frac{1}{2}\sum_k\left[\frac{r_k^2}{\epsilon_k^2 + \sigma_b^2 + \sigma_s^2} + \ln(\epsilon_k^2 + \sigma_b^2 + \sigma_s^2) + \ln 2\pi\right].
\end{equation*}
Here $\sigma_s$ and $\sigma_b$ appear only as $\sigma_s^2 + \sigma_b^2$, so they are exactly degenerate.

\paragraph{One epoch.}
$A = 1/d$ and $B = r/d$, so Eq.~\eqref{eq:chi2} gives $r^2/(d + \sigma_s^2)$ and Eq.~\eqref{eq:logdet} gives $\ln(d + \sigma_s^2)$:
a 1D gaussian with variance $\epsilon^2 + \sigma_b^2 + \sigma_s^2$.

\paragraph{A missing epoch.}
The code gives it $\epsilon_k = 10^4$ km/s (\texttt{missing\_epoch\_err}), so $1/d_k \approx 10^{-8}$.
That epoch drops out of $A$ and $B$, and the star reduces to the one-epoch case.
Its $\ln d_k \approx \ln 10^8$ is a constant that doesn't depend on the parameters
(up to $\sigma_b^2/10^8$).

\section{Another route: integrating out the offset}
The same result follows without any matrix algebra, and this route makes the role of ``shared'' explicit.
Given $s$, the epochs are independent: $v_k \mid s \sim \mathcal{N}(\mu + s, d_k)$.
$s$ is unknown, so integrate over its distribution:
\begin{equation}
  p(\bv) = \int_{-\infty}^{\infty} \left[\prod_{k=1}^{n}\mathcal{N}(v_k \mid \mu + s, d_k)\right]\mathcal{N}(s\mid 0, \sigma_s^2)\, ds .
\end{equation}
One integration variable appears in every epoch's factor; that is what ``shared'' means.
In the independent model each epoch has its own $s_k$, so the integral splits into $n$ separate ones,
each of which just adds $\sigma_s^2$ to $d_k$.

The exponent is
\begin{equation*}
  -\frac{1}{2}\left[\sum_k\frac{(r_k - s)^2}{d_k} + \frac{s^2}{\sigma_s^2}\right]
  = -\frac{1}{2}\left[\sum_k\frac{r_k^2}{d_k} - 2sB + s^2 P\right], \qquad P \equiv A + \frac{1}{\sigma_s^2}.
\end{equation*}
Complete the square, $s^2P - 2sB = P(s - B/P)^2 - B^2/P$, and use
$\int e^{-P(s-c)^2/2}\,ds = \sqrt{2\pi/P}$.
The normalizing prefactors are $\prod_k(2\pi d_k)^{-1/2}$ and $(2\pi\sigma_s^2)^{-1/2}$. Then
\begin{equation*}
  p(\bv) = \prod_k (2\pi d_k)^{-1/2}\;\frac{1}{\sqrt{1 + \sigma_s^2 A}}\;
  \exp\!\left(-\frac{1}{2}\left[\sum_k \frac{r_k^2}{d_k} - \frac{\sigma_s^2B^2}{1 + \sigma_s^2 A}\right]\right),
\end{equation*}
using $\sigma_s^2 P = 1 + \sigma_s^2 A$ and $B^2/P = \sigma_s^2 B^2/(1 + \sigma_s^2 A)$.
Taking $-\ln$ gives Eq.~\eqref{eq:result}.

\section{Caveat}
The algebra is exact for the model as stated.
But the model assumes every star's binary motion is $\mathcal{N}(0, \sigma_b^2)$ with one shared $\sigma_b$.
The real binary motion is heavy-tailed: most stars are singles with $b_k = 0$, and a few binaries move by tens of km/s.
That mis-specification is why $\sigma_s \to 0$ in the \texttt{fit\_shared} scan
(see the readme section on \texttt{single\_component\_comparisons\_mockobs.py}).
A per-star mixture (binary with probability $f_b$) would replace $p(\bv)$ by
$f_b\, p(\bv \mid \sigma_b) + (1 - f_b)\, p(\bv \mid \sigma_b = 0)$.
Each term is Eq.~\eqref{eq:result}.

\end{document}
